\subsection{Die Raumladungszone im pn-Übergang}
% Alt: Aufgabe 2
\Aufgabe{
  \label{Aufg2}
  Untersuchen Sie die Eigenschaften und die Bildung der Raumladungszone in pn-Übergängen von Halbleitern. 
  Beantworten Sie die folgenden Aspekte in Ihrer Ausarbeitung:
  \begin{itemize}
    \item Was ist die Raumladungszone und wie entsteht diese im pn-Übergang?
    \item Wie ist der Verlauf des elektrischen Feldes im Halbleiter?
    \item Welche inneren Vorgänge wirken auf die freien Ladungsträger?
    \item Wie verändert sich die Raumladungszone bei einer angelegten Spannung in Durchlass- und Sperrrichtung?
          
  \end{itemize}
}


\Loesung{
  Die RLZ ist eine schmale 
  Region um den pn-Übergang herum, in der positive Ionen aus der n-Schicht und negative Ionen aus der p-Schicht verbleiben, 
  nachdem die Diffusion abgeschlossen ist. In dieser Zone gibt es keine freien Ladungsträger, da die positiven und negativen Ladungen sich 
  gegenseitig neutralisieren. Dadurch entsteht ein elektrisches Feld ($\vec{E}$), das sowohl einen Driftstrom erzeugt als auch die Diffusion von weiteren 
  Ladungsträgern unterdrückt. Die Raumladungszone wirkt wie eine Sperrschicht und verhindert den Stromfluss in Sperrrichtung.
  Siehe: Abschnitt \ref{sec:pn-Übergang}
}